\BourbakiExercisesSection

*1. 在仿射空间 \(\mathbf{E}\) （在域 \(\mathbf{K}\) 上）中，四元组 \((a, b, c, d)\) （ \(\mathbf{E}\) 的点的四元组）称为平行四边形，如果 \(b - a = c - d\) ，这时 \((a, d, c, b)\) 也是一个平行四边形。证明若 \(\mathbf{K}\) 的特征 \(\neq 2\) ，则有序对 \((a, c)\) 和 \((b, d)\) 的中点相等；当 \(\mathbf{K}\) 的特征为 \(2?_{*}\) 。

*2. 设 E 是特征 \(\neq 2\) 的域上的一个仿射空间，而 \(a, b, c, d\) 是 E 的任意四个点。证明若 \(x, y, z, t\) 分别是有序对 \((a, b), (b, c), (c, d), (d, a)\) 的中点，则四元组 \((x, y, z, t)\) 是一个平行四边形（习题 1）。*

*3. 设 K 是特征 \(\neq 2\) 的交换域，E 是 K 上的一个仿射平面，而 \(a, b, c, d\) 是 E 的四个点，其中任意三点都不共线。设 \(\mathbf{D}_{xy}\) 表示过 E 的两个不同点 \(x, y\) 的直线。假设直线 \(\mathbf{D}_{ab}\) 和 \(\mathbf{D}_{cd}\) 有一个公共点 \(e\) ，且 \(\mathbf{D}_{ad}\) 和 \(\mathbf{D}_{bc}\) 有一个公共点 \(f\) 。证明三个有序对 \((a, c), (b, d), (e, f)\) 的中点位于同一条直线上。当 \(\mathbf{D}_{ab}\) 和 \(\mathbf{D}_{cd}\) 平行或 \(\mathbf{D}_{ad}\) 和 \(\mathbf{D}_{bc}\) 平行时，这一性质变成什么？考虑 K 是 3 个元素的域的情形。*

*4. 设 K 是特征异于 2 和 3 的域，E 是 K 上的一个仿射空间， \(a, b, c\) 是 E 的不在同一条直线上的三个点，而 \(a', b', c'\) 分别是有序对 \((b, c)\) , \((c, a)\) 和 \((a, b)\) 的中点。证明（用习题 3 的记号）直线 \(\mathbf{D}_{aa'}\) , \(\mathbf{D}_{bb'}\) 和 \(\mathbf{D}_{cc'}\) 都过三个点 \(a, b, c\) 的重心。当 K 的特征为 2 或特征为 3 时，这一性质变成什么？推广到 \(n\) 个仿射无关点的系统的情形。*

\begin{enumerate}
\def\labelenumi{\arabic{enumi}.}
\setcounter{enumi}{4}
\item
  设 E 是域 K 上的一个仿射空间，且 K 至少含三个元素。为使 E 的非空子集 V 是一个线性簇，必要且充分条件是：对 V 中每对 \((x,y)\) 不同的点 x 和 y，过 x 和 y 的直线 \(D_{xy}\) 整个包含在 V 中。若 K 恰有两个元素，则为使 V 是一个线性簇，必要且充分条件是 V 的任意三点的重心属于 V。
\item
  \begin{enumerate}
  \def\labelenumii{(\alph{enumii})}
  \tightlist
  \item
    设 E 是域 K 上的一个 \(\geqslant2\) 维仿射空间。为使 E 到自身的仿射映射 u 把 E 的每条直线变为平行直线，必要且充分条件是：与 u 相伴的线性映射 v 是一个位似 \(t\mapsto\gamma t\) （比例 \(\gamma\neq0\) 属于 K 的中心）。若 \(\gamma=1\) ，u 是一个平移；证明若 \(\gamma\neq1\) ，存在且仅存在一个点 \(a\in E\) 使得 \(u(a)=a\) 。若取 a 作为 E 的原点，则 u 在 E 上由此确定的向量空间结构下等同于一个中心位似；u 称为仿射空间 E 的以 a 为中心、以 \(\gamma\) 为比例的中心位似。
  \end{enumerate}
\end{enumerate}

\begin{enumerate}
\def\labelenumi{(\alph{enumi})}
\setcounter{enumi}{1}
\item
  设 \(u_1, u_2\) 是 \(\mathbf{E}\) 到 \(\mathbf{E}\) 的两个仿射映射，其中每一个要么是平移，要么是中心位似。证明 \(u_1 \circ u_2\) 是 \(\mathbf{E}\) 的一个平移或一个中心位似；若 \(u_1, u_2\) 和 \(u_1 \circ u_2\) 这三个都是中心位似，证明它们的中心位于同一条直线上。当 \(u_1\) 和 \(u_2\) 是中心位似而 \(u_1 \circ u_2\) 是平移时，能说些什么？
\item
  证明平移集与中心位似之集构成仿射群的一个正规子群 H，并且 H/T 同构于 K 的中心的乘法群；证明 H 仅当 H = T 时才可能是交换的，换言之，仅当 K 的中心恰含两个元素时。
\end{enumerate}

¶ 7. 设 E（相应地 E'）是域 K（至少含 3 个元素）上的一个有限维 \(n \geqslant 2\) 仿射空间（相应地，域 K' 上的仿射空间），并设 u 是 E 到 E' 的一个单射映射，它把 E 中同一条直线上任意三个点映到同一条直线上的三个点，并且使得由 \(u(E)\) 在 E' 中生成的仿射线性簇等于 E'。

\begin{enumerate}
\def\labelenumi{(\alph{enumi})}
\item
  证明 \(u\) 把 \(E\) 的每个仿射无关点系统映为一个仿射无关点系统（利用习题 5）。
\item
  设 \(D_{1}, D_{2}\) 是 E 中两条平行直线，而 \(D_{1}^{\prime}, D_{2}^{\prime}\) 是 \(E^{\prime}\) 中分别包含 \(u(\mathbf{D}_{1})\) 和 \(u(\mathbf{D}_{2})\) 的直线。证明 \(D_{1}^{\prime}\) 和 \(D_{2}^{\prime}\) 位于同一平面内；若进一步 u 是满射，证明 \(D_{1}^{\prime}\) 和 \(D_{2}^{\prime}\) 平行（在相反情形，证明 E 中会有 3 个不在同一直线上的点，它们在 u 下的像在同一直线上）（参见习题 17）。
\item
  此后假设：若 \(D_{1}\) 和 \(D_{2}\) 是 E 中的平行直线，则 \(E'\) 中分别包含 \(u(\mathbf{D}_{1})\) 和 \(u(\mathbf{D}_{2})\) 的直线是平行的。证明若取 E 中的原点 a 及原点 \(a' = u(a)\) （在 \(E'\) 中），则存在一个同构 \(\sigma\) （把 K 映到一个子域 \(K_{1}\) （属于 \(K'\) ）上），使得若把 E 视为 K 上的向量空间，且把 \(E'\) 作为 \(K_{1}, u\) 上的向量空间，则 u 是一个单射半线性映射（相对于 \(\sigma\) ，由 E 到 \(E'\) ）（§ 1，第 13 节）。（先考虑 n = 2 的情形；给定 E 的一个基 \((e_{1}, e_{2})\) ，证明对 K 的任意两个元素 \(\alpha, \beta\) ，点 \((\alpha + \beta)e_{1}\) 和 \((\alpha\beta)e_{1}\) （属于 E）可以从点 0、 \(e_{1}, e_{2}, \alpha e_{1}, \beta e_{1}\) 出发，通过作给定直线的平行线及给定直线的交点而构造出来；由此推出 \(u(\lambda e_{1}) = \lambda^{\sigma}u(e_{1})\) ，其中 \(\sigma\) 是 K 到 \(K'\) 的一个子域上的同构，然后证明也有 \(u(\lambda e_{2}) = \lambda^{\sigma}u(e_{2})\) （通过考虑连接点 \(\lambda e_{1}\) 和 \(\lambda e_{2}\) 的直线）。最后由此过渡到 n 为任意的情形，对 n 作归纳论证。）若 u 是双射，证明 \(K_{1} = K'\) 。
\item
  把 (c) 的结果推广到 K 是含 2 个元素的域的情形，同时假设 u 也把 E 的每个仿射无关点系统映为 E' 的仿射无关点系统。
\end{enumerate}

\begin{enumerate}
\def\labelenumi{\arabic{enumi}.}
\setcounter{enumi}{7}
\tightlist
\item
  设 E 是域 K 上的左仿射空间，T 为其平移空间。
\end{enumerate}

\begin{enumerate}
\def\labelenumi{(\alph{enumi})}
\tightlist
\item
  为使映射 \(f\) （由 \(\mathbf{E}\) 到一个左向量空间 \(\mathbf{L}\) （在 \(\mathbf{K}\) 上））是仿射的，必要且充分条件是
\end{enumerate}

\[
\begin{array}{r l} f (\mathbf {t} + x) - f (x) & = f (\mathbf {t} + y) - f (y) \\ f (\lambda \mathbf {t} + x) - f (x) & = \lambda f (\mathbf {t} + x) - f (x)) \end{array}
\]

对所有 \(\lambda\in\mathbf{K},\mathbf{t}\in\mathbf{T},x,y\) （在 E 中）。设 \(x\mapsto[x]\) 是 E 到由 E 的元素的形式线性组合所成的向量空间 \(\mathrm{K}_{s}^{(\mathbb{E})}\) 中的典范单射，并设 N 是 \(\mathrm{K}_{s}^{(\mathbb{E})}\) 的由元素

\[
[ \mathbf {t} + x ] - [ x ] - [ \mathbf {t} + y ] + [ y ]
\]

\[
[ \lambda \mathbf {t} + x ] - [ x ] - \lambda [ \mathbf {t} + x ] + \lambda [ x ]
\]

对 \(\lambda \in \mathbf{K}\) , \(\mathbf{t} \in \mathbf{T}\) , \(x, y\) （在 \(\mathbf{E}\) 中）；最后，令 \(\mathbf{V}\) 为 \(\mathbf{K}_s^{(\mathbf{E})}\) 模 \(\mathbf{N}\) 的商空间，并令 \(\phi\) 为 \(\mathbf{E}\) 到 \(\mathbf{V}\) 的映射，它把每个 \(x \in \mathbf{E}\) 对应到 \([x]\) 模 \(\mathbf{N}\) 的类。则 \(\phi\) 是一个仿射映射，并且对每个仿射映射 \(f\) （由 \(\mathbf{E}\) 到一个左向量空间 \(\mathbf{L}\) （在 \(\mathbf{K}\) 上）），存在且仅存在一个线性映射 \(g\) （由 \(\mathbf{V}\) 到 \(\mathbf{L}\) ）使得 \(f = g \circ \phi\) 。

\begin{enumerate}
\def\labelenumi{(\alph{enumi})}
\setcounter{enumi}{1}
\tightlist
\item
  设 \(\phi_0: \mathbf{T} \to \mathbf{V}\) 是与 \(\phi\) 相伴的线性映射，它使得 \(\phi(\mathbf{t} + x) - \phi(x) = \phi_0(\mathbf{t})\) 。证明 \(\phi\) （从而 \(\phi_0\) 也是）是单射（考虑映射 \(x \mapsto x - a\) （由 \(\mathbf{E}\) 到 \(\mathbf{T}\) ，对某个 \(a \in \mathbf{E}\) ））；对每个族 \((\lambda_t)_{t \in I}\) （ \(\mathbf{K}\) 的元素，具有有限支撑）与每个族 \((x_t)_{t \in I}\) （ \(\mathbf{E}\) 的元素），
\end{enumerate}

\[
\sum_ {i} \lambda_ {i} \phi (x _ {i}) = \phi_ {0} \left(\sum_ {i} \lambda_ {i} x _ {i}\right) \quad \text { if } \quad \sum_ {i} \lambda_ {i} = 0
\]

\[
\sum_ {i} \lambda_ {i} \phi (x _ {i}) = \mu \phi_ {0} \left(\sum_ {i} \mu^ {- 1} \lambda_ {i} x _ {i}\right) \quad \text { if } \quad \sum_ {i} \lambda_ {i} = \mu \neq 0.
\]

由此推出 \(\phi_{0}(T)\) 是 V 中过 0 的一个超平面，而 \(\phi(E)\) 是与 \(\phi_{0}(T)\) 平行的一个超平面。

\begin{enumerate}
\def\labelenumi{\arabic{enumi}.}
\setcounter{enumi}{8}
\tightlist
\item
  设 \(\mathbf{K}\) 是一个含 \(q\) 个元素的有限域，而 \(\mathbf{V}\) 是一个 \(n\) 维向量空间（在 \(\mathbf{K}\) 上）。
\end{enumerate}

\begin{enumerate}
\def\labelenumi{(\alph{enumi})}
\tightlist
\item
  证明序列 \((x_{1}, x_{2}, \ldots, x_{m})\) （ \(m \leqslant n\) 个 \(\mathbf{V}\) 的向量，构成自由系统）所成的集的基数等于
\end{enumerate}

\[
(q ^ {n} - 1) (q ^ {n} - q) \dots (q ^ {n} - q ^ {m - 1})
\]

（对 \(n\) 作归纳论证）。

\begin{enumerate}
\def\labelenumi{(\alph{enumi})}
\setcounter{enumi}{1}
\tightlist
\item
  由 (a) 推出，\(m\) 维线性簇（在一个 \(n\) 维射影空间中，空间在域 \(K\) 上）所成的集的基数等于
\end{enumerate}

\[
\frac {(q ^ {n + 1} - 1) (q ^ {n + 1} - q) \dots (q ^ {n + 1} - q ^ {m})}{(q ^ {m + 1} - 1) (q ^ {m + 1} - q) \dots (q ^ {m + 1} - q ^ {m})}.
\]

\begin{enumerate}
\def\labelenumi{\arabic{enumi}.}
\setcounter{enumi}{9}
\item
  在一个 \(n\) 维射影空间 \(\mathbf{P}(\mathbf{V})\) （在域 \(\mathbf{K}\) 上）中，射影基是一个集合 \(\mathbf{S}\) （由 \(n + 2\) 个点组成），其中任意两点都构成一个射影自由系。若 \(\mathbf{S} = (a_i)_{0 \leqslant i \leqslant n + 1}\) 和 \(\mathbf{S}' = (a_i')_{0 \leqslant i \leqslant n + 1}\) 是 \(\mathbf{P}(\mathbf{V})\) 的任意两个射影基，证明存在一个变换 \(f \in \mathbf{PGL}(\mathbf{V})\) 使得 \(f(a_i) = a_i'\) （对 \(0 \leqslant i \leqslant n + 1\) ）。为使这个变换唯一，必要且充分条件是 \(\mathbf{K}\) 可交换。（把它归结为 \(a_i' = a_i\) （对所有 \(i\) ）成立的情形，并注意总可以写成（用第 6 节的记号） \(a_i = \pi(b_i)\) ，其中 \((b_i)_{1 \leqslant i \leqslant n + 1}\) 是 \(\mathbf{V}\) 的一个基，而 \(b_0 = b_1 + b_2 + \cdots + b_{n+1}\) 。） *给出一个例子，其中 \(\mathbf{K}\) 是四元数域，并且存在点的一个无限集 \(\mathbf{T}\) （属于 \(\mathbf{P}(\mathbf{V})\) ），其中任意 \(n + 1\) 个点都构成射影自由系，又存在一个变换 \(f \in \mathbf{PGL}(\mathbf{V})\) ，它不同于恒等变换，且保持 \(\mathbf{T}\) 的所有点不变。*
\item
  设 V 是域 K 上的一个二维向量空间，而 a, b, c, d 是射影直线 \(\mathbf{P}(\mathbf{V})\) 上的四个不同点。四元组 \((a, b, c, d)\) 的交比（记作 \(\begin{bmatrix} a & b \\ d & c \end{bmatrix}\) ）是使下述条件成立的元素 \(\xi \in \mathbf{K}\) 所成的集：存在两个向量 \(u, v\) （在 \(\mathbf{V}\) 中），对它们（用第 6 节的记号）有 \(a = \pi(u)\) , \(b = \pi(v)\) , \(c = \pi(u + v)\) , \(d = \pi(u + \xi v)\) 。这个定义立即推广到具有射影直线结构的集合（第 11 节）的任意四个不同点的四元组。
\end{enumerate}

\begin{enumerate}
\def\labelenumi{(\alph{enumi})}
\item
  证明 \(\begin{bmatrix} a & b \\ d & c \end{bmatrix}\) 是一个元素 \(\neq 1\) （属于乘法群 \(\mathbf{K}^*\) ）的共轭元素的集；反之，当 \(a, b, c\) 是 \(\mathbf{P}(V)\) 的不同点而 \(\rho\) 是一个元素 \(\neq 1\) （属于 \(\mathbf{K}^*\) ）的共轭元素的集时，存在一个点 \(d \in \mathbf{P}(V)\) 使得 \(\begin{bmatrix} a & b \\ d & c \end{bmatrix} = \rho\) 。为使 \(d\) 唯一，必要且充分条件是 \(\rho\) 由单个点组成。
\item
  证明
\end{enumerate}

\[
\left[ \begin{array}{c c} a & b \\ c & d \end{array} \right] = \left[ \begin{array}{c c} b & a \\ d & c \end{array} \right] = \left[ \begin{array}{c c} a & b \\ d & c \end{array} \right] ^ {- 1}
\]

和

\[
\left[ \begin{array}{c c} d & a \\ c & b \end{array} \right] = 1 - \left[ \begin{array}{c c} a & b \\ d & c \end{array} \right]
\]

（用 \(\rho^{-1}\) （相应地， \(1 - \rho\) ）记共轭元素的集 \(\lambda \xi^{-1} \lambda^{-1} = (\lambda \xi \lambda^{-1})^{-1}\) （相应地， \(1 - \lambda \xi \lambda^{-1} = \lambda (1 - \xi) \lambda^{-1}\) ），其中 \(\xi\) 是 \(\rho\) 的一个元素）。

\begin{enumerate}
\def\labelenumi{(\alph{enumi})}
\setcounter{enumi}{2}
\tightlist
\item
  设 \((a, b, c, d)\) , \((a', b', c', d')\) 是 \(\mathbf{P}(\mathbf{V})\) 的两组由不同点组成的四元组。为使存在 V 到自身的一个双射半线性映射，使得在过渡到商时得到的 \(\mathbf{P}(\mathbf{V})\) 到自身的双射映射 f 满足条件 \(f(a) = a'\) , \(f(b) = b'\) , \(f(c) = c'\) , \(f(d) = d'\) ，必要且充分条件是存在 K 的一个自同构 \(\sigma\) 使得
\end{enumerate}

\[
\left[ \begin{array}{c c} a ^ {\prime} & b ^ {\prime} \\ d ^ {\prime} & c ^ {\prime} \end{array} \right] = \left[ \begin{array}{c c} a & b \\ d & c \end{array} \right] ^ {\sigma}.
\]

为使存在射影群 PGL(V) 的一个变换 f 满足上述条件，必要且充分条件是

\[
\left[ \begin{array}{c c} a ^ {\prime} & b ^ {\prime} \\ d ^ {\prime} & c ^ {\prime} \end{array} \right] = \left[ \begin{array}{c c} a & b \\ d & c \end{array} \right].
\]

\begin{enumerate}
\def\labelenumi{\arabic{enumi}.}
\setcounter{enumi}{11}
\tightlist
\item
  设 \(\mathbf{P}(\mathbf{V})\) 是一个有限维数 \(n\) 的（左）射影空间（在域 \(\mathbf{K}\) 上）。证明在 \(\mathbf{P}(\mathbf{V})\) 的射影超平面所成的集上存在一个 \(n\) 维（右）射影空间结构（在 \(\mathbf{K}\) 上），它典范同构于 \(\mathbf{P}(\mathbf{V}^*)\) 的结构（ \(\mathbf{V}^*\) 是 \(\mathbf{V}\) 的对偶）。若 \(\mathbf{M}\) 是维数为 \(r < n\) 的一个线性簇（在 \(\mathbf{P}(\mathbf{V})\) 中），推出一个 \((n - r - 1)\) 维射影空间结构（在包含 \(\mathbf{M}\) 的射影超平面所成的集上）。特别地，若 \(\mathbf{M}\) 的维数为 \(n - 2\) ，可以定义交比 \(\begin{bmatrix} \mathrm{H}_1 & \mathrm{H}_2 \\ \mathrm{H}_4 & \mathrm{H}_3 \end{bmatrix}\) —— 包含 M 的不同超平面的四元组 \((\mathrm{H}_1, \mathrm{H}_2, \mathrm{H}_3, \mathrm{H}_4)\) 的交比。证明若 \(\mathbf{D} \subset \mathbf{P}(\mathbf{V})\) 是一条与 M 不相交的直线，而 \(a_i\) 是 D 与 \(\mathrm{H}_i\) 的交点 ( \(1 \leqslant i \leqslant 4\) )，则
\end{enumerate}

\[
\left[ \begin{array}{c c} a _ {1} & a _ {2} \\ a _ {4} & a _ {3} \end{array} \right] = \left[ \begin{array}{c c} \mathrm{H} _ {1} & \mathrm{H} _ {2} \\ \mathrm{H} _ {4} & \mathrm{H} _ {3} \end{array} \right].
\]

*13. 在射影平面 \(\mathbf{P}(\mathbf{V})\) （在域 \(\mathbf{K}\) 上，特征 \(\neq 2\) ）中，设 \(a, b, c, d\) 是构成一个射影基的四个点（习题 10）；记 \(\mathbf{D}_{xy}\) 为过两个不同点 \(x, y\) （属于 \(\mathbf{P}(\mathbf{V})\) ）的直线，设 \(e, f, g\) 分别是直线 \(\mathbf{D}_{ab}\) 与 \(\mathbf{D}_{cd}\) 、 \(\mathbf{D}_{ac}\) 与 \(\mathbf{D}_{bd}\) 、 \(\mathbf{D}_{ad}\) 与 \(\mathbf{D}_{bc}\) 的交点；令 \(h\) 为 \(\mathbf{D}_{bc}\) 与 \(\mathbf{D}_{ef}\) 的交点；证明 \(\begin{bmatrix} b & c \\ h & g \end{bmatrix} = \{-1\}\) （“完全四边形定理”；把它归结为 \(\mathbf{D}_{ad}\) 是一个仿射平面的无穷远直线的情形）。当 \(\mathbf{K}\) 的特征为 \(2?_{*}\)

\begin{enumerate}
\def\labelenumi{\arabic{enumi}.}
\setcounter{enumi}{13}
\item
  在射影平面 \(\mathbf{P}(\mathrm{V})\) （在至少含三个元素的域 \(\mathbf{K}\) 上）中，设 \(\mathbf{D}, \mathbf{D}'\) 为两条不同的直线。为使对任意不同的点 \(a, b, c\) 属于 \(\mathbf{D}, a', b', c'\) 属于 \(\mathbf{D}'\) ，交点 \(r\) （ \(\mathbf{D}_{ab'}\) 与 \(\mathbf{D}_{ba'}\) 的交点）、 \(q\) （ \(\mathbf{D}_{ac'}\) 与 \(\mathbf{D}_{ca'}\) 的交点）、 \(p\) （ \(\mathbf{D}_{bc'}\) 与 \(\mathbf{D}_{cb'}\) 的交点）位于同一条直线上，必要且充分条件是 \(\mathbf{K}\) 可交换（“帕普斯定理”；把它归结为 \(q\) 和 \(r\) 位于一个仿射平面的无穷远直线上的情形）。把这一定理应用于 \(\mathbf{P}(\mathrm{V})\) 的直线的射影空间（习题 12）。
\item
  在域 K（至少含三个元素）上的射影平面中，设 \(s, a, b, c, a', b', c'\) 为七个不同的点，使得 \(\{s, a, b, c\}\) 和 \(\{s, a', b', c'\}\) 是射影基（习题 10），且直线 \(\mathbf{D}_{sa}, \mathbf{D}_{sb}, \mathbf{D}_{sc}\) 分别过 \(a', b', c'\) 。证明交点 \(r\) （ \(\mathbf{D}_{ab}\) 与 \(\mathbf{D}_{a'b'}\) 的交点）、 \(p\) （ \(\mathbf{D}_{bc}\) 与 \(\mathbf{D}_{b'c'}\) 的交点）、 \(q\) （ \(\mathbf{D}_{ca}\) 与 \(\mathbf{D}_{c'a'}\) 的交点）位于同一条直线上（“德萨格定理”；方法类似于习题 14）。
\item
  \begin{enumerate}
  \def\labelenumii{(\alph{enumii})}
  \tightlist
  \item
    设 \(\mathbf{E} = \mathbf{P}(\mathbf{V})\) 和 \(\mathbf{E}' = \mathbf{P}(\mathbf{V}')\) 是两个同维数 \(n \geqslant 2\) 的射影空间（分别在域 \(\mathbf{K}\) , \(\mathbf{K}'\) 上），并设 \(u\) 是 \(\mathbf{E}\) 到 \(\mathbf{E}'\) 上的一个双射映射，它把同一条直线上任意三个点映到同一条直线上的三个点。证明存在一个同构 \(\sigma\) （把 \(\mathbf{K}\) 映到 \(\mathbf{K}'\) 上）以及一个双射半线性映射 \(v\) （ \(\mathbf{V}\) 到 \(\mathbf{V}'\) 上，相对于 \(\sigma\) ），使得 \(u\) 是在过渡到商时由 \(v\) 得到的映射（“射影几何基本定理”；利用习题 7）。进一步假设 \(\mathbf{V}' = \mathbf{V}\) 且 \(\mathbf{K}\) 是交换的；为使 \(u\) 是一个射影映射，必要且充分条件是 \(\begin{bmatrix} u(a) & u(b) \\ u(d) & u(c) \end{bmatrix} = \begin{bmatrix} a & b \\ d & c \end{bmatrix}\) 对每个四元组 \((a, b, c, d)\) （由 \(\mathbf{P}(\mathbf{V})\) 中同一条直线上的不同点组成）也成立。
  \end{enumerate}
\end{enumerate}

\begin{enumerate}
\def\labelenumi{(\alph{enumi})}
\setcounter{enumi}{1}
\tightlist
\item
  设 \(p\) 是一个使 \(1 \leqslant p \leqslant n - 1\) 的整数。证明当在 \(u\) 下每个 \(p\) 维射影线性簇的像包含在一个 \(p\) 维射影线性簇中时，(a) 的第一个结论成立。
\end{enumerate}

\begin{enumerate}
\def\labelenumi{\arabic{enumi}.}
\setcounter{enumi}{16}
\tightlist
\item
  设 V 是域 K 上有限维数 n 的一个向量空间， \((e_{i})_{1\leq i\leq n}\) 是 V 的一组基， \(K'\) 是 K 的一个子域，而 \(V'\) 是 \(K'\) 上的 n 维向量空间（由 \(e_{i}\) 生成）。给出 \(V'\) 到 V 的一个单射映射的例子，它把仿射空间 \(V'\) 中同一条直线上任意三个点映到仿射空间 V 中同一条直线上的三个点，但不一定把两条平行直线映到包含在两条平行直线中的集合。（把 V 典范地嵌入与其相伴的射影空间 E 中，并考虑 E 到自身的一个射影变换 u，使得无穷远超平面在 u 下的原像不同于该超平面且不含 \(V'\) 的任何点；例如可取 K 为无限而 \(K'\) 为有限。）
\end{enumerate}

¶ 18. 设 \(\mathbf{E} = \mathbf{P}(\mathbf{V})\) 是域 \(\mathbf{K}\) 上的一个射影平面，而 \(u\) 是 \(\mathbf{E}\) 到自身的一个双射映射，它在过渡到商时由一个双射半线性映射 \(v\) （ \(\mathbf{V}\) 到自身，相对于一个自同构 \(\sigma\) —— \(\mathbf{K}\) 的自同构）导出。

\begin{enumerate}
\def\labelenumi{(\alph{enumi})}
\item
  证明以下四个性质等价：(α) 对所有 \(x \in E\) , \(x, u(x)\) 和 \(u^{2}(x)\) 位于同一条直线上；（β） E 的每条直线都含有一个在 u 下不变的点；（γ）过 E 的每个点都有一条在 u 下不变的直线；（δ）对 E 的每条直线 D，三条直线 D、 \(u(\mathbf{D})\) 和 \(u^{2}(\mathbf{D})\) 有一个公共点。（首先证明 (α) 与 (γ) 等价；利用对偶性（习题 12）推出 (β) 与 (δ) 等价；最后证明 (γ) 蕴含 (β)，并利用对偶性推出 (β) 蕴含 (γ)。）
\item
  设 u 具有 (a) 中所述的性质。证明：若在 E 中存在一条在 u 下不变的直线 D，且它仅含一个在 u 下不变的点，则 u 在过渡到商时由 V 的一个平延 v 导出（§ 10，习题 11）。（证明每条在 u 下不变的直线都含 a；通过考虑一条不过 a 的直线，证明存在一条过 a 且至少含两个在 u 下不变点的直线 \(D_{0}\) ；由此推出 \(D_{0}\) 的所有点都必然在 u 下不变。）
\item
  假设 u 具有 (a) 中所述的性质。证明：若在 E 中存在一条在 u 下不变的直线 E，且它仅含两个在 u 下不变的点，则 u 在过渡到商时由 V 的一个伸缩 v 导出（ \(\S\) 10，习题 11）。（若 a、b 是 D 中在 u 下不变的两个点，证明每条在 u 下不变的直线都过 a 或过 b；然后注意至少存在另外两个不同于 a、b 且在 u 下不变的点 c、d，因而直线 \(D_{cd}\) 过 a 或 b；通过证明 \(D_{cd}\) 的所有点都在 u 下不变来得出结论。）
\item
  假设 u 具有 (a) 中所述的性质，并且 E 中每条在 u 下不变的直线都至少含三个在 u 下不变的不同点；这时 E 中存在一个每点都在 u 下不变的射影基（习题 10）；由此推出存在 V 的一组基 \((e_{i})_{1\leq i\leq3}\) ，使得 u 在过渡到商时由一个半线性映射 \(v\) （ \(\mathbf{V}\) 到自身）导出，它满足 \(v(e_i) = e_i\) （对 \(1 \leqslant i \leqslant 3\) ）。 \(\mathbf{E}\) 中在 \(u\) 下不变的点的集这时是射影平面 \(\mathbf{P}(\mathbf{V}_0)\) ，其中 \(\mathbf{V}_0\) 是 \(\mathbf{K}_0\) （ \(\sigma\) 的不变量所成的域）上由 \(e_1, e_2, e_3\) 生成的向量空间。
\item
  此后假设 \(u\) 满足 (a) 和 (d) 的条件，且 \(u\) 和 \(u^2\) 都不是恒等映射。证明存在 \(\gamma \in \mathbf{K}\) 使得 \(\gamma^\sigma = \gamma\) 且
\end{enumerate}

\[
(\xi^ {\sigma} - \xi) ^ {\sigma} = \gamma (\xi^ {\sigma} - \xi)\tag{1}
\]

对所有 \(\xi \in \mathbf{K}\) 成立（利用条件 \((\alpha)\) （属于 \((a)\) ）以及存在 \(\zeta \in \mathbf{K}\) 使得 \(\zeta^{\sigma} \neq \zeta\) ）。证明 \(\gamma \neq -1\) ，并且对所有 \(\xi \in \mathbf{K}\) （使 \(\xi^{\sigma} \neq \xi\) ），

\[
(1 + \gamma) \xi^ {\sigma} \gamma = \gamma \xi (1 + \gamma)\tag{2}
\]

（把（1）应用于把 \(\xi\) 换成 \(\xi^2\) 的情形）；把（2）推广到所有 \(\xi \in \mathbf{K}\) ，办法是注意 \(\xi = \eta - \zeta\) （其中 \(\eta^\sigma \neq \eta\) 且 \(\xi^\sigma \neq \zeta\) ）。由此得出结论 \(\gamma \neq 1\) ，然后由（1）和（2）推出

\[
\xi^ {\sigma} = (1 + \gamma) \xi (1 + \gamma) ^ {- 1}\tag{3}
\]

对所有 \(\xi \in K\) 成立，并且 \(\gamma + \gamma^{-1}\) 属于 \(K\) 的中心。证明逆命题。*给出一个例子，其中 \(\gamma\) 不属于 \(K\) 的中心，但 \(\gamma + \gamma^{-1}\) 属于该中心。*

*19. 设 K 为一个交换域， \(\tilde{K}\) 为通过向 K 添加一个无穷远点（第 9 节）所得的射影域，而 \(f\) 和 \(g\) 为 K(X) 的两个元素。证明，记 \(h(\mathbf{X}) = f(g(\mathbf{X}))\) ，若 \(\tilde{f}, \tilde{g}, \tilde{h}\) 是 \(f, g, h\) 到 \(\tilde{\mathbf{K}}\) 的典范扩张，则 \(\tilde{h} = \tilde{f} \circ \tilde{g}_{\cdot *}\)

